\documentclass[../DoAn.tex]{subfiles}
\begin{document}

Phụ lục này tổng hợp các test case kiểm thử chức năng (black-box testing) của hệ thống, bao gồm 28 test case bao phủ các module xác thực, tải lên CV, soạn thảo và xuất CV. Các test case được thực hiện thủ công theo kịch bản; với các component React quan trọng, khóa luận bổ sung unit test bằng Jest và Testing Library. Tất cả 28 test case đều đạt; chi tiết từng module được trình bày dưới đây.

\section{Kiểm thử module xác thực}

Bảng~\ref{tab:test-auth} trình bày các trường hợp kiểm thử cho module xác thực, bao phủ đăng ký, đăng nhập, OAuth và kiểm soát truy cập.

\begin{table}[H]
\centering
\caption{Kết quả kiểm thử module xác thực}
\label{tab:test-auth}
\begin{tabular}{|c|p{5cm}|p{5cm}|c|}
\hline
\textbf{TC} & \textbf{Mô tả} & \textbf{Kết quả mong đợi} & \textbf{Trạng thái} \\
\hline
TC01 & Đăng ký với email hợp lệ & Tạo tài khoản thành công & Đạt \\
\hline
TC02 & Đăng ký với email đã tồn tại & Thông báo lỗi email trùng & Đạt \\
\hline
TC03 & Đăng ký với mật khẩu yếu & Thông báo yêu cầu mật khẩu mạnh & Đạt \\
\hline
TC04 & Đăng nhập bằng email/mật khẩu đúng & Đăng nhập thành công, chuyển hướng & Đạt \\
\hline
TC05 & Đăng nhập bằng username (admin alias) & Đăng nhập thành công & Đạt \\
\hline
TC06 & Đăng nhập với mật khẩu sai & Thông báo lỗi xác thực & Đạt \\
\hline
TC07 & Đăng nhập bằng Google OAuth & Chuyển hướng Google, callback thành công & Đạt \\
\hline
TC08 & Đăng nhập bằng LinkedIn OAuth & Chuyển hướng LinkedIn, callback thành công & Đạt \\
\hline
TC09 & Truy cập trang bảo vệ khi chưa đăng nhập & Chuyển hướng đến trang đăng nhập & Đạt \\
\hline
TC10 & Truy cập CV editor khi là khách & Cho phép truy cập (guest session) & Đạt \\
\hline
\end{tabular}
\end{table}

\section{Kiểm thử module tải lên và phân tích CV}

Bảng~\ref{tab:test-upload} liệt kê các trường hợp kiểm thử cho module tải lên CV.

\begin{table}[H]
\centering
\caption{Kết quả kiểm thử module tải lên và phân tích CV}
\label{tab:test-upload}
\begin{tabular}{|c|p{5cm}|p{5cm}|c|}
\hline
\textbf{TC} & \textbf{Mô tả} & \textbf{Kết quả mong đợi} & \textbf{Trạng thái} \\
\hline
TC12 & Upload file PDF hợp lệ (<10MB) & Trích xuất văn bản và phân tích thành công & Đạt \\
\hline
TC13 & Upload file DOCX hợp lệ & Trích xuất và phân tích thành công & Đạt \\
\hline
TC14 & Upload file không hỗ trợ (\texttt{.jpg}) & Thông báo lỗi định dạng & Đạt \\
\hline
TC15 & Upload file quá 10MB & Thông báo vượt kích thước cho phép & Đạt \\
\hline
TC16 & Upload CV tiếng Việt & Nhận diện ngôn ngữ VI, phân tích đúng & Đạt \\
\hline
TC17 & Upload CV tiếng Anh & Nhận diện ngôn ngữ EN, phân tích đúng & Đạt \\
\hline
TC18 & Upload file không phải CV & \texttt{possibility\_score} thấp (<30) & Đạt \\
\hline
TC19 & Phân tích khi OpenAI API lỗi & Chuyển sang fallback parser & Đạt \\
\hline
\end{tabular}
\end{table}

\section{Kiểm thử trình soạn thảo CV}

Bảng~\ref{tab:test-editor} liệt kê các trường hợp kiểm thử cho trình soạn thảo CV.

\begin{table}[H]
\centering
\caption{Kết quả kiểm thử trình soạn thảo CV}
\label{tab:test-editor}
\begin{tabular}{|c|p{5cm}|p{5cm}|c|}
\hline
\textbf{TC} & \textbf{Mô tả} & \textbf{Kết quả mong đợi} & \textbf{Trạng thái} \\
\hline
TC20 & Chỉnh sửa thông tin liên hệ & Cập nhật và hiển thị trên preview & Đạt \\
\hline
TC21 & Thêm mục kinh nghiệm mới & Thêm thành công, hiển thị trên preview & Đạt \\
\hline
TC22 & Xóa mục kinh nghiệm & Xóa thành công với xác nhận & Đạt \\
\hline
TC23 & Kéo thả sắp xếp thứ tự & Thứ tự cập nhật trên cả editor và preview & Đạt \\
\hline
TC24 & Thêm/xóa kỹ năng dạng tag & Tag thêm/xóa đúng & Đạt \\
\hline
TC25 & Auto-save sau chỉnh sửa & Dữ liệu lưu tự động sau 2 giây & Đạt \\
\hline
TC26 & Preview cập nhật real-time & Preview phản ánh ngay khi sửa & Đạt \\
\hline
TC27 & Thêm section tùy chỉnh & Section mới hiển thị đúng & Đạt \\
\hline
\end{tabular}
\end{table}

\section{Kiểm thử module xuất CV}

Bảng~\ref{tab:test-export} liệt kê các trường hợp kiểm thử cho module xuất CV.

\begin{table}[H]
\centering
\caption{Kết quả kiểm thử module xuất CV}
\label{tab:test-export}
\begin{tabular}{|c|p{5cm}|p{5cm}|c|}
\hline
\textbf{TC} & \textbf{Mô tả} & \textbf{Kết quả mong đợi} & \textbf{Trạng thái} \\
\hline
TC28 & Xuất PDF & Mở print dialog, layout đúng & Đạt \\
\hline
TC29 & Xuất DOCX & Tải file \texttt{.docx}, mở được bằng Word & Đạt \\
\hline
\end{tabular}
\end{table}

\section{Đánh giá Hybrid Parser so với Full-LLM Parser}

Để đánh giá hiệu quả của phương pháp hybrid parser (đóng góp \#1 trong §1.5), thí nghiệm so sánh được thực hiện giữa hai phương pháp: (i) Full-LLM parser (sử dụng hoàn toàn GPT cho toàn bộ CV), và (ii) Hybrid parser (regex + LLM cho phần kinh nghiệm). Bộ dữ liệu thử nghiệm gồm 20 CV mẫu (10 tiếng Việt, 10 tiếng Anh) với độ dài và cấu trúc đa dạng.

\begin{table}[H]
\centering
\caption{So sánh hiệu năng Hybrid Parser vs Full-LLM Parser}
\label{tab:hybrid-vs-fullllm}
\begin{tabular}{|l|c|c|}
\hline
\textbf{Tiêu chí} & \textbf{Full-LLM} & \textbf{Hybrid} \\
\hline
Thời gian trung bình (giây) & 8.5 & 4.2 \\
\hline
Token sử dụng trung bình & 3200 & 1400 \\
\hline
Chi phí ước tính / CV (USD) & 0.0048 & 0.0021 \\
\hline
Độ chính xác thông tin liên hệ & 95\% & 93\% \\
\hline
Độ chính xác kinh nghiệm & 90\% & 88\% \\
\hline
Độ chính xác kỹ năng & 88\% & 85\% \\
\hline
Tỷ lệ giảm token & -- & 56\% \\
\hline
\end{tabular}
\end{table}

Kết quả trong Bảng~\ref{tab:hybrid-vs-fullllm} cho thấy phương pháp hybrid giảm được 56\% lượng token API sử dụng và giảm hơn một nửa thời gian xử lý so với phương pháp Full-LLM. Mức giảm nhẹ về độ chính xác (2--3\%) ở các trường thông tin liên hệ và kỹ năng là chấp nhận được, do phần lớn thông tin này có cấu trúc tương đối cố định và regex hoạt động đủ tốt. Đối với phần kinh nghiệm làm việc, mức chênh lệch 2\% (90\% vs 88\%) chủ yếu đến từ các trường hợp CV có cấu trúc không chuẩn, nơi LLM có lợi thế hơn trong việc hiểu ngữ cảnh.

\section{So sánh với các giải pháp hiện có}

Bảng~\ref{tab:compare-existing} so sánh chi tiết hệ thống với bốn giải pháp phổ biến trên thị trường.

\begin{table}[H]
\centering
\caption{So sánh chi tiết hệ thống với các giải pháp hiện có}
\label{tab:compare-existing}
\small
\begin{tabular}{|p{2.8cm}|c|c|c|c|c|}
\hline
\textbf{Tiêu chí} & \textbf{Canva} & \textbf{Zety} & \textbf{Sovren} & \textbf{LinkedIn} & \textbf{Khoá luận} \\
\hline
Upload CV có sẵn & Không & Không & Có & Không & Có \\
\hline
Phân tích bằng AI/LLM & Không & Hạn chế & NLP truyền thống & Không & LLM (GPT) \\
\hline
Trình soạn thảo tương tác & Template-based & Form-based & Không & Profile-based & Section-based + DnD \\
\hline
Xem trước real-time & Có & Có & Không & Không & Có \\
\hline
Xuất PDF & Có & Có & Không & Hạn chế & Có \\
\hline
Xuất DOCX & Không & Có & Không & Không & Có \\
\hline
Hỗ trợ tiếng Việt & Hạn chế & Không & Hạn chế & Hạn chế & Đầy đủ \\
\hline
Đa phương thức xác thực & Email & Email & API key & LinkedIn & Email + Google + LinkedIn \\
\hline
Phiên khách (Guest) & Không & Không & Không & Không & Có \\
\hline
Auto-save & Có & Có & -- & Có & Có \\
\hline
Mã nguồn mở & Không & Không & Không & Không & Có \\
\hline
Chi phí & Freemium & Trả phí & Trả phí cao & Miễn phí & Miễn phí \\
\hline
\end{tabular}
\end{table}

Hệ thống do khoá luận xây dựng có lợi thế rõ rệt ở các điểm: (i) kết hợp được cả tải lên CV, phân tích bằng trí tuệ nhân tạo, và soạn thảo tương tác trong cùng một hệ thống, (ii) hỗ trợ tiếng Việt đầy đủ, (iii) cho phép phiên khách sử dụng mà không cần đăng ký, và (iv) mã nguồn mở cho phép cộng đồng đóng góp và tùy chỉnh.

\end{document}
